\documentclass[11pt]{article}
\usepackage[T1]{fontenc}
\usepackage[a4paper,margin=24mm]{geometry}
\usepackage{amsmath,amssymb}
\usepackage[hidelinks]{hyperref}
\usepackage{microtype}

\setlength{\parindent}{0pt}
\setlength{\parskip}{0.58em}

\newcommand{\G}{\mathbb G}
\newcommand{\X}{\mathbb X}
\newcommand{\wideG}{\widehat{\mathbb G}}
\newcommand{\wideX}{\widehat{\mathbb X}}

\title{Coamenability versus Relative Amenability for Coideals:\\
A Scoped Later Answer}
\author{Literature status for arXiv:2003.04384}
\date{13 August 2026}

\begin{document}
\maketitle

\begin{center}
\fbox{\parbox{0.91\textwidth}{
\textbf{Status: literature already answered in two broad regimes; the
unrestricted question remains open.}
Anderson--Sackaney and Khosravi, arXiv:2311.14587, prove the desired
coamenability/relative-amenability equivalence when the ambient discrete
quantum group is unimodular or exact.  In arbitrary ambient groups they prove
one implication and obtain a full equivalence after strengthening relative
amenability to equivariant injectivity.
}}
\end{center}

\section*{The source question}

Question~1.3 of Anderson--Sackaney's source paper asks the following.  Let
$\G$ be a discrete quantum group and let $\widehat{\mathbb H}$ be a closed
quantum subgroup of $\wideG$.  Is
\[
 L^\infty(\wideG/\widehat{\mathbb H})\ \text{coamenable}
 \quad\Longleftrightarrow\quad
 \ell^\infty(\mathbb H)\ \text{relatively amenable}?
\]
The next sentence gives the compact-quasi-subgroup version: for an idempotent
state $\omega$, is the compact coideal $N_\omega$ coamenable exactly when its
codual coideal $\widetilde N_{P_\omega}$ is relatively amenable?  The source
explicitly says that it only makes progress on this version~\cite{source}.
The question and extension occur on source PDF page~4.

\section*{The later theorem}

The supporting paper uses the paired notation
\[
 \wideX\backslash\wideG
 \qquad\text{and}\qquad
 \ell^\infty(\X)
\]
for a compact-quasi-subgroup coideal and its codual.  This is the same
coideal--codual formulation as the source's $N_\omega$ and
$\widetilde N_{P_\omega}$; quotient-side symbols depend on the chosen
left/right convention.

Its two main implications are:
\begin{align*}
 \wideX\backslash\wideG\text{ coamenable}
 &\Longrightarrow \ell^\infty(\X)\text{ is $\G$-injective}
 &&\text{(Theorem~3.2, PDF page~12)},\\
 \ell^\infty(\X)\text{ is $\G$-injective}
 &\Longrightarrow \wideX\backslash\wideG\text{ coamenable}
 &&\text{(Theorem~6.5, PDF page~25)}.
\end{align*}
Thus, without any unimodularity or exactness assumption,
\[
 \boxed{
 \wideX\backslash\wideG\text{ is coamenable}
 \iff
 \ell^\infty(\X)\text{ is $\G$-injective}.}
\]
Since $\G$-injectivity implies relative amenability, the forward implication
asked for in the source follows for every compact quasi-subgroup.  These
statements are author-explicit: the abstract and introduction of the
supporting paper state this coideal/codual equivalence directly~\cite{answer}.

\section*{When the requested equivalence is complete}

Remark~6.6 on supporting PDF page~27 proves the missing implication
\[
 \ell^\infty(\X)\text{ relatively amenable}
 \Longrightarrow
 \wideX\backslash\wideG\text{ coamenable}
\]
under either of the following assumptions:
\begin{enumerate}
 \item $\G$ is unimodular;
 \item $C(\wideX\backslash\wideG)$ is exact.
\end{enumerate}
In particular, the second condition holds whenever $\G$ is exact, because
$C(\wideX\backslash\wideG)$ is a $C^*$-subalgebra of $C(\wideG)$.
Combining Remark~6.6 with Theorem~3.2 therefore yields the exact equivalence
asked for in the compact-quasi-subgroup extension whenever $\G$ is
unimodular or exact.  Coideals of quotient type are included in this class,
so the corresponding closed-quantum-subgroup question has the same scoped
affirmative answer.

The proof mechanisms also explain the two hypotheses.  In the unimodular
case, Theorem~6.2 upgrades relative amenability to $\G$-injectivity.  In the
exact case, the proof of Theorem~6.5 can use a nuclear inclusion supplied by
injectivity of the von Neumann coideal in place of nuclearity of the reduced
$C^*$-coideal.

\section*{What is not settled}

The later paper does not prove that relative amenability implies
$\G$-injectivity for an arbitrary non-unimodular, non-exact discrete quantum
group, nor does it prove the converse implication in the source question
without one of the above hypotheses.  Its unrestricted exact replacement is
$\G$-injectivity, which is strictly stronger at the level of the definitions.
Accordingly, this packet records a substantial scoped literature answer, not
a full resolution of Question~1.3 in complete generality.

\section*{Verification and provenance}

The source question and compact-quasi-subgroup extension were checked on
source PDF page~4.  In the supporting paper, the summary equivalence was
checked in the abstract (page~1) and introduction (page~2), the universal
forward implication in Theorem~3.2 (page~12), the injectivity converse in
Theorem~6.5 (page~25), and the unimodular/exact upgrade in Remark~6.6
(page~27).  The supporting paper is coauthored by the source author and cites
the source for the coideal framework.  No novelty claim is made here.

\begin{thebibliography}{9}
\bibitem{source}
B. Anderson--Sackaney,
\emph{On amenable and coamenable coideals},
J. Funct. Anal. \textbf{286} (2024), no.~10, 110384;
arXiv:2003.04384.

\bibitem{answer}
B. Anderson--Sackaney and F. Khosravi,
\emph{Topological boundaries of representations and coideals},
arXiv:2311.14587v2 (2024).
\end{thebibliography}

\end{document}
